% Add your packages here
% \usepackage[backend=biber]{biblatex}
% \addbibresource{references.bib}
\usepackage[english]{babel}
\usepackage[autostyle]{csquotes}
\usepackage[backref=page]{hyperref}
\usepackage[usenames,dvipsnames]{xcolor}
\usepackage{tikz}
%\usetikzlibrary{external}

\usepackage{amssymb,amsmath,amsfonts,amsthm}
\usepackage{xspace}
\usepackage{bm}
\usepackage{nicefrac}
\usepackage[usenames,dvipsnames]{xcolor} %Options expand the named colours
\usepackage{framed}
%\usepackage{duckuments}
%\usepackage[algo2e]{algorithm2e}
%\usepackage{palatino} %preferred font2
%\usepackage{tgpagella} %preferred font1
%\usepackage{mathptmx} %preferred font3 = times new roman style
%\usepackage{mathpazo}
%\usepackage{dsfont}
\usepackage{indentfirst}
\usepackage{setspace}
\usepackage{fancyhdr}
\usepackage{colortbl}
%\usepackage[table, xcdraw]{xcolor}
\pagestyle{plain}
\fancyhf{}
\renewcommand{\headrulewidth}{0pt} 
\usepackage[numbers, sort&compress]{natbib}
%\usepackage[square,comma,compress]{natbib}
\bibliographystyle{unsrtnat}

\usepackage[capitalise]{cleveref}
\usepackage[toc,page]{appendix}
\usepackage{mathtools}
\usepackage{booktabs}
\usepackage{array}
\usepackage{colortbl}
\usepackage{graphicx}
\usepackage{siunitx}
\usepackage{ifxetex}

\usepackage{soul}
\usepackage{everysel}
\usepackage{verbatim}
\usepackage{tabularx}


%\usepackage{microtype}
\usepackage[strict]{changepage}

\newlength{\alphabet}
\settowidth{\alphabet}{\normalfont abcdefghijklmnopqrstuvwxyz}

\usepackage{geometry}

\makeatletter
%\geometry{hmarginratio={1:1}}
\geometry{hmarginratio={1:1}, margin=1in} % fallback

\@ifclassloaded{book}
{
  \if@twoside
  \geometry{hmarginratio={2:3}}
  \else
  \geometry{hmarginratio={1:1}}
  \fi
}{}

\@ifclassloaded{article}
{
  \if@twoside
  \geometry{hmarginratio={4:2}}
  \else
  \geometry{hmarginratio={1:1}}
  \fi
}{}
\makeatletter

%http://texblog.net/latex-archive/plaintex/full-justification-with-typewriter-font
\EverySelectfont{%
  \fontdimen2\font=0.4em% interword space
  \fontdimen3\font=0.2em% interword stretch
  \fontdimen4\font=0.1em% interword shrink
  \fontdimen7\font=0.1em% extra space
  \hyphenchar\font=`\-% to allow hyphenation
}

%% new types of table columns
\newcolumntype{L}[1]{>{\raggedright\let\newline\\\arraybackslash\hspace{0pt}}m{#1}}
\newcolumntype{C}[1]{>{\centering\let\newline\\\arraybackslash\hspace{0pt}}m{#1}}
\newcolumntype{R}[1]{>{\raggedleft\let\newline\\\arraybackslash\hspace{0pt}}m{#1}}
\newcolumntype{J}[1]{>{\let\newline\\\arraybackslash\hspace{0pt}}m{#1}}
\newcolumntype{j}[1]{>{\let\newline\\\arraybackslash\hspace{0pt}}b{#1}}

%% spacing in table rows
\setstretch{1.0}
\renewcommand{\arraystretch}{1.2}

\usepackage{subcaption}
\setlength{\captionmargin}{0.025\textwidth}

\setlength{\parskip}{0.5\baselineskip}%
%\setlength{\parindent}{0pt}%

\definecolor{darkgray}{rgb}{0.2,0.2,0.2}

%% set up hyperref
\hypersetup{%
%  backref=page,
%  bookmarks=true,
  pdftoolbar=false,
  pdfmenubar=true,
  pdfstartview={FitH},
  pdftitle={},
  pdfauthor={Yuheng Qiu},
  pdfsubject={},
  pdfkeywords={}{},
  colorlinks=true,
  citecolor=blue,
  linkcolor=blue,
  urlcolor=blue
}

\newtheorem{mydef}{Definition}
\newtheorem{prop}{Proposition}[section]

%% common math
\DeclareMathOperator*{\diag}{diag}
\DeclareMathOperator*{\trace}{trace}
\DeclareMathOperator*{\rank}{rank}
\DeclareMathOperator*{\argmax}{arg\,max}
\DeclareMathOperator*{\argmin}{arg\,min}
\DeclareRobustCommand*{\median}{median}
\newcommand{\smallplus}{\raisebox{.3ex}{$\mkern1mu\scriptstyle+$}}
\newcommand{\Argmin}[1]{\underset{#1}{\argmin}}
\newcommand{\Argmax}[1]{\underset{#1}{\argmax}}
\newcommand{\norm}[1]{\lVert#1\rVert}

%\newcommand{\T}{\ensuremath{^\top}}
% http://tex.stackexchange.com/questions/36817/how-to-typeset-the-subscript-of-a-matrix
%\ifxetex
%\renewcommand{\T}{^{\mathstrut\scriptscriptstyle{\top}}} % ooh, fancy!
%\else
\newcommand{\T}{^{\mathstrut\scriptscriptstyle{\top}}} % ooh, fancy!
%\fi
\newcommand{\mc}[1]{\ensuremath{\mathcal{#1}}}   % caligraphic math typeface
\newcommand{\mb}[1]{\ensuremath{\mathbf{#1}}}    % bold math typeface
\newcommand{\mbb}[1]{\ensuremath{\mathbb{#1}}}
\newcommand{\mcb}[1]{\ensuremath{\bm{\mathcal{#1}}}}
\newcommand{\LC}[1]{\ensuremath{\,#1\mathcal{L}}}
\newcommand{\MC}[1]{\ensuremath{\,#1\mathcal{M}}}
\newcommand{\nth}[1]{\ensuremath{#1^\text{th}}}
\newcommand{\sk}[1]{\left[#1\right]_\times}  % skew symmetric matrix
\newcommand{\set}[3]{\ensuremath{\left\{ #1 \right\}_{#2}^{#3}}}

\makeatletter
\DeclareRobustCommand\onedot{\futurelet\@let@token\@onedot}
\def\@onedot{\ifx\@let@token.\else.\null\fi\xspace}

\def\eg{\emph{e.g}\onedot} \def\Eg{\emph{E.g}\onedot}
\def\ie{\emph{i.e}\onedot} \def\Ie{\emph{I.e}\onedot}
\def\cf{\emph{c.f}\onedot} \def\Cf{\emph{C.f}\onedot}
\def\etc{\emph{etc}\onedot} \def\vs{\emph{vs}\onedot}
\def\wrt{w.r.t\onedot}
\def\etal{\emph{et al}\onedot}
\def\aka{\emph{a.k.a}\onedot}
\def\dof{d.o.f\onedot}

\newcommand{\tform}[2]{{\,^{\mathtt{#2}}\mb{T}\,_{\mathtt{#1}}}}

\usepackage{ifdraft}
\ifdraft{%
  \usepackage{draftwatermark}
  \SetWatermarkText{DRAFT}
  \SetWatermarkAngle{55}
  \SetWatermarkScale{6.0}
  \SetWatermarkLightness{0.85}
  \SetWatermarkFontSize{12 pt}
}

\usepackage{etoolbox}

\apptocmd{\thebibliography}{\csname phantomsection\endcsname\addcontentsline{toc}{chapter}{\bibname}}{}{}

\newlength\fwidth
\newlength\fheight

% Redefine \vec with a better one
\usepackage[T1]{fontenc}  % For correct {}s in \texttt
%\usepackage[b]{esvect}    % For \vv

% --- Macro \xvec
\makeatletter
\newlength\xvec@height%
\newlength\xvec@depth%
\newlength\xvec@width%
\newcommand{\xvec}[2][]{%
  \ifmmode%
    \settoheight{\xvec@height}{$#2$}%
    \settodepth{\xvec@depth}{$#2$}%
    \settowidth{\xvec@width}{$#2$}%
  \else%
    \settoheight{\xvec@height}{#2}%
    \settodepth{\xvec@depth}{#2}%
    \settowidth{\xvec@width}{#2}%
  \fi%
  \def\xvec@arg{#1}%
  \def\xvec@dd{:}%
  \def\xvec@d{.}%
  \raisebox{.2ex}{\raisebox{\xvec@height}{\rlap{%
    \kern.05em%  (Because left edge of drawing is at .05em)
    \begin{tikzpicture}[scale=1]
    \pgfsetroundcap
    \draw (.05em,0)--(\xvec@width-.05em,0);
    \draw (\xvec@width-.05em,0)--(\xvec@width-.15em, .075em);
    \draw (\xvec@width-.05em,0)--(\xvec@width-.15em,-.075em);
    \ifx\xvec@arg\xvec@d%
      \fill(\xvec@width*.45,.5ex) circle (.5pt);%
    \else\ifx\xvec@arg\xvec@dd%
      \fill(\xvec@width*.30,.5ex) circle (.5pt);%
      \fill(\xvec@width*.65,.5ex) circle (.5pt);%
    \fi\fi%
    \end{tikzpicture}%
  }}}%
  #2%
}
\makeatother

% Define some commenting commands
\usepackage{xargs}                      % Use more than one optional parameter in a new commands
%\usepackage[colorinlistoftodos,prependcaption,textsize=tiny]{todonotes}
%\newcommandx{\unsure}[2][1=]{\todo[linecolor=red,backgroundcolor=red!25,bordercolor=red,#1]{#2}}
%\newcommandx{\todopicture}[2][1=]{\todo[linecolor=blue,backgroundcolor=blue!25,bordercolor=blue,#1]{#2}}
%\newcommandx{\info}[2][1=]{\todo[linecolor=OliveGreen,backgroundcolor=OliveGreen!25,bordercolor=OliveGreen,#1]{#2}}
%\newcommandx{\improvement}[2][1=]{\todo[linecolor=Plum,backgroundcolor=Plum!25,bordercolor=Plum,#1]{#2}}
%\newcommandx{\thiswillnotshow}[2][1=]{\todo[disable,#1]{#2}}

% Change the figure caption size
\captionsetup[figure]{font=footnotesize,labelfont={footnotesize, bf}}

\usepackage{longtable}
%\usepackage{arydshln}
\usepackage{enumitem}

\usepackage{tabularx}
\usepackage{transparent}
\usepackage{multicol}
\usepackage{multirow}
\usepackage{makecell}
\usepackage{colortbl}

% Define a multiline cell in the tables
\newenvironment{multilinecell}% environment name
{ \renewcommand{\arraystretch}{1.5}%
  \begin{tabular}[t]{@{} p{\linewidth} @{}}%
}%
{\end{tabular}}
% Define an environment for challenges
\newcounter{challenge}[section]
\newenvironment{challenge}[2][]
{ \refstepcounter{challenge}\par\medskip
 \noindent \textbf{Challenge~\thechallenge.~~#2:} \par}
{\medskip}


\newcommand*\matrice[1]{\begin{bmatrix}#1\end{bmatrix}}

% Allow matrix environment to accept the vertical spacing:
% Example \begin{bmatrix}[1.5]
\makeatletter
\renewcommand*\env@matrix[1][\arraystretch]{%
  \edef\arraystretch{#1}%
  \hskip -\arraycolsep
  \let\@ifnextchar\new@ifnextchar
  \array{*\c@MaxMatrixCols c}}
\makeatother

% Force subfigures to reference like 12(a)
\captionsetup[subfigure]{subrefformat=simple,labelformat=simple}
    \renewcommand\thesubfigure{(\alph{subfigure})}


% ------------------------------------------------------------------------
% Math symbol types
% \let\stdvec\vec                                         % redefine the standard \vec as \stdvec

% \newcommand{\FrameSymbol}{\mc{F}}                       % F symbol for frame
% \newcommand{\SetSymbol}{\mc{S}}                       % F symbol for frame
% \renewcommand{\vec}[1]{\xvec[]{#1}}                     % vector
% \newcommand{\dvec}[1]{\xvec[.]{#1}}                     % dotted vector
% \newcommand{\ddvec}[1]{\xvec[:]{#1}}                    % double-dotted vector
% \newcommand{\frm}[1]{\mc{#1}}                           % reference frame letter
% \newcommand{\frmsub}[2][]{\mc{#2}_{#1}}                 % reference frame letter with subscript
% \newcommand{\Frm}[1]{{{\FrameSymbol}^{\frm{#1}}}}       % reference frame with its F
% \newcommand{\Frmsub}[2][]{{\FrameSymbol}^{\frmsub[#1]{#2}}} % reference frame with its F
% \newcommand{\mat}[1]{{\mathbf{#1}}}                     % matrix
% \newcommand{\numset}[3]{{\mathbb{#1}}^{{#2} \times {#3}}} % number set (e.g., R^3x1)
% \newcommand{\point}[1]{{#1}}
% \newcommand{\axis}[1]{\hat{#1}}
% \newcommand{\unit}[1]{\left[ #1 \right]}
% \newcommand{\matcomp}[2]{{\left[ #1 \right]}_{#2}}
% \newcommand{\vecelem}[2]{{_{#2}{#1}}}
% \newcommand{\Project}[2]{{_{#2}{#1}}}
% \newcommand{\Plane}[2]{{#1 #2}}

% \newcommand{\Set}[1]{\SetSymbol_\mc{#1}}       % Font for sets
% \newcommand{\Seti}[2]{\SetSymbol_{\mc{#1}_#2}}       % Font for sets

% ---------------------------------------------------------------------------
\definecolor{commentcolor}{gray}{0.5}
\usepackage{algorithm}
\usepackage{algpseudocode}
\renewcommand{\algorithmicforall}{\textbf{for each}}
\algnewcommand{\LineComment}[1]{\State \textcolor{commentcolor}{\(\triangleright\) #1}}
\algnewcommand{\To}{\textbf{to}}
\algnewcommand{\Break}{\textbf{break}}
\algnewcommand{\Continue}{\textbf{continue}}
\algnewcommand{\IIf}[1]{\State\algorithmicif\ #1\ \algorithmicthen}
\algnewcommand{\EndIIf}{\unskip}
\algnewcommand{\var}[1]{\textit{#1}}
\algnewcommand{\func}[1]{\textsc{#1}}
\renewcommand{\algorithmicrequire}{\textbf{Input:}}
\renewcommand{\algorithmicensure}{\textbf{Output:}}


% ------------------------------------------------------------------------
% Kinematics symbols
\newcommand{\FI}{\Frm{I}}
\newcommand{\FB}{\Frm{B}}
\newcommand{\FS}{\Frm{S}}
\newcommand{\FC}{\Frm{C}}
\newcommand{\FE}{\Frm{E}}
\newcommand{\FR}[1]{\Frmsub[#1]{R}}

\newcommand{\OI}{\point{O}_{\frm{I}}}
\newcommand{\OB}{\point{O}_{\frm{B}}}
\newcommand{\OR}[1]{\point{O}_{\frmsub[#1]{R}}}
\newcommand{\ORB}[1]{_{\frm{B}}\point{O}_{\frmsub[#1]{R}}}

\newcommand{\XI}{\axis{X}_{\frm{I}}}
\newcommand{\YI}{\axis{Y}_{\frm{I}}}
\newcommand{\ZI}{\axis{Z}_{\frm{I}}}
\newcommand{\XB}{\axis{X}_{\frm{B}}}
\newcommand{\YB}{\axis{Y}_{\frm{B}}}
\newcommand{\ZB}{\axis{Z}_{\frm{B}}}
\newcommand{\XR}[1]{\axis{X}_{\frmsub[#1]{R}}}
\newcommand{\YR}[1]{\axis{Y}_{\frmsub[#1]{R}}}
\newcommand{\ZR}[1]{\axis{Z}_{\frmsub[#1]{R}}}
\newcommand{\XS}{\axis{X}_{\frm{S}}}
\newcommand{\YS}{\axis{Y}_{\frm{S}}}
\newcommand{\ZS}{\axis{Z}_{\frm{S}}}
\newcommand{\XC}{\axis{X}_{\frm{C}}}
\newcommand{\YC}{\axis{Y}_{\frm{C}}}
\newcommand{\ZC}{\axis{Z}_{\frm{C}}}
\newcommand{\XE}{\axis{X}_{\frm{E}}}
\newcommand{\YE}{\axis{Y}_{\frm{E}}}
\newcommand{\ZE}{\axis{Z}_{\frm{E}}}
\newcommand{\XpB}{\axis{X'}_{\frm{B}}}
\newcommand{\YpB}{\axis{Y'}_{\frm{B}}}
\newcommand{\ZpB}{\axis{Z'}_{\frm{B}}}
\newcommand{\XF}{\axis{X}_{\frm{F}}}
\newcommand{\YF}{\axis{Y}_{\frm{F}}}
\newcommand{\ZF}{\axis{Z}_{\frm{F}}}

\newcommand{\iI}{\axis{\mat{i}}_{\frm{I}}}
\newcommand{\jI}{\axis{\mat{j}}_{\frm{I}}}
\newcommand{\kI}{\axis{\mat{k}}_{\frm{I}}}
\newcommand{\iIv}{\axis{\vec{i}}_{\frm{I}}}
\newcommand{\jIv}{\axis{\vec{j}}_{\frm{I}}}
\newcommand{\kIv}{\axis{\vec{k}}_{\frm{I}}}
\newcommand{\iB}{\axis{\mat{i}}_{\frm{B}}}
\newcommand{\jB}{\axis{\mat{j}}_{\frm{B}}}
\newcommand{\kB}{\axis{\mat{k}}_{\frm{B}}}
\newcommand{\iBv}{\axis{\vec{i}}_{\frm{B}}}
\newcommand{\jBv}{\axis{\vec{j}}_{\frm{B}}}
\newcommand{\kBv}{\axis{\vec{k}}_{\frm{B}}}
\newcommand{\ipB}{\axis{\mat{i'}}_{\frm{B}}}
\newcommand{\jpB}{\axis{\mat{j'}}_{\frm{B}}}
\newcommand{\kpB}{\axis{\mat{k'}}_{\frm{B}}}
\newcommand{\ipBv}{\axis{\vec{i'}}_{\frm{B}}}
\newcommand{\jpBv}{\axis{\vec{j'}}_{\frm{B}}}
\newcommand{\kpBv}{\axis{\vec{k'}}_{\frm{B}}}
\newcommand{\iS}{\axis{\mat{i}}_{\frm{S}}}
\newcommand{\jS}{\axis{\mat{j}}_{\frm{S}}}
\newcommand{\kS}{\axis{\mat{k}}_{\frm{S}}}
\newcommand{\iSv}{\axis{\vec{i}}_{\frm{S}}}
\newcommand{\jSv}{\axis{\vec{j}}_{\frm{S}}}
\newcommand{\kSv}{\axis{\vec{k}}_{\frm{S}}}
%\newcommand{\iFv}{\axis{\vec{i}}_{\frm{F}}}
%\newcommand{\jFv}{\axis{\vec{j}}_{\frm{F}}}
%\newcommand{\kFv}{\axis{\vec{k}}_{\frm{F}}}
\newcommand{\iCv}{\axis{\vec{i}}_{\frm{C}}}
\newcommand{\jCv}{\axis{\vec{j}}_{\frm{C}}}
\newcommand{\kCv}{\axis{\vec{k}}_{\frm{C}}}

\newcommand{\raxis}{\vec{r}}

\newcommand{\RBI}{\mat{R}_{\frm{BI}}}
\newcommand{\RIB}{\mat{R}_{\frm{IB}}}
\newcommand{\RIS}{\mat{R}_{\frm{IS}}}
\newcommand{\RSI}{\mat{R}_{\frm{SI}}}
\newcommand{\RIC}{\mat{R}_{\frm{IC}}}
\newcommand{\RCI}{\mat{R}_{\frm{CI}}}
\newcommand{\RIE}{\mat{R}_{\frm{IE}}}
\newcommand{\REI}{\mat{R}_{\frm{EI}}}
\newcommand{\REC}{\mat{R}_{\frm{EC}}}
\newcommand{\RCE}{\mat{R}_{\frm{CE}}}
\newcommand{\RRB}[1]{\mat{R}_{\frmsub[#1]{R}\frm{B}}}
\newcommand{\RBR}[1]{\mat{R}_{\frm{B}\frmsub[#1]{R}}}
\newcommand{\InertB}{\mat{I}_{\frm{B}}}

\newcommand{\phix}[1]{\phi_{x{#1}}}
\newcommand{\phiy}[1]{\phi_{y{#1}}}

\DeclareMathOperator{\cosine}{c}
\DeclareMathOperator{\sine}{s}
\newcommand{\ctheta}{\cosine{\theta}}
\newcommand{\stheta}{\sine{\theta}}
\newcommand{\cphi}{\cosine{\phi}}
\newcommand{\sphi}{\sine{\phi}}
\newcommand{\cpsi}{\cosine{\psi}}
\newcommand{\spsi}{\sine{\psi}}
\newcommand{\cmu}[1]{\cosine{\mu_{#1}}}
\newcommand{\smu}[1]{\sine{\mu_{#1}}}
\newcommand{\cphix}[1]{\cosine{\phix{#1}}}
\newcommand{\sphix}[1]{\sine{\phix{#1}}}
\newcommand{\cphiy}[1]{\cosine{\phiy{#1}}}
\newcommand{\sphiy}[1]{\sine{\phiy{#1}}}

\DeclareMathOperator{\Span}{Span}
\DeclareMathOperator{\Conv}{Conv}

% Attitude generation symbols

\newcommand{\spt}{_{sp}}
\newcommand{\des}{_{des}}
\newcommand{\ver}{_{ver}}
\newcommand{\hover}{_{hov}}
\newcommand{\hor}{_{hor}}
\newcommand{\lat}{_{lat}}
\newcommand{\nor}{_{nor}}

\newcommand{\Flmax}{F_{lmax}}
\newcommand{\Flmaxnew}{{F'}_{lmax}}

\newcommand{\Fdes}{\vec{F}\des}
\newcommand{\FdesI}{\vec{F}\des^\frm{I}}
\newcommand{\Fdesnew}{\vec{F'}\des}
\newcommand{\FdesInew}{\vec{F}\des^{\prime\,\frm{I}}}
\newcommand{\FdesCf}{\stdvec{F}\des^{f\,\frm{C}}}
\newcommand{\FdesIf}{\stdvec{F}\des^{f\,\frm{I}}}
\newcommand{\FdesIp}{\stdvec{F}\des^{p\,\frm{I}}}
\newcommand{\MdesCm}{\stdvec{M}\des^{m\,\frm{C}}}

\newcommand{\Fwind}{\vec{F}_{wind}}
\newcommand{\Fgrav}{\vec{F}_{grav}}
\newcommand{\Fext}{\vec{F}_{ext}}

\newcommand{\Fspt}{\vec{F}\spt}
\newcommand{\FsptB}{\vec{F}\spt^\frm{B}}

\newcommand{\Fver}{\vec{F}\ver}
\newcommand{\FverB}{\vec{F}\ver^{\frm{B}}}

\newcommand{\Flat}{\vec{F}\lat}
\newcommand{\FlatI}{\vec{F}\lat^{\frm{I}}}
\newcommand{\Fnor}{\vec{F}\nor}
\newcommand{\FnorI}{\vec{F}\nor^{\frm{I}}}

\newcommand{\Fhor}{\vec{F}\hor}
\newcommand{\FhorB}{\vec{F}\hor^{\frm{B}}}
\newcommand{\Fhornew}{\vec{F'}\hor}
\newcommand{\FhorBnew}{\vec{F}\hor^{\prime\,\frm{B}}}

\newcommand{\Fsptnew}{\vec{F'}\spt}
\newcommand{\FsptBnew}{\vec{F}\spt^{\prime\,\frm{B}}}
\newcommand{\FsptBnewnew}{\vec{F}\spt^{''\,\frm{B}}}

\newcommand{\PlaneB}{{\Plane{\XB}{\YB}}}
\newcommand{\PlaneI}{{\Plane{\XI}{\YI}}}
\newcommand{\PlaneS}{{\Plane{\XS}{\YS}}}
\newcommand{\PlaneC}{{\Plane{\XC}{\YC}}}

\newcommand{\ProjB}[1]{\Project{#1}{\PlaneB}}
\newcommand{\ProjI}[1]{\Project{#1}{\PlaneI}}

\newcommand{\Fhover}{F\hover}
\newcommand{\FhovB}{\vec{F}\hover^\frm{B}}
\newcommand{\FhovI}{\vec{F}\hover^\frm{I}}

\newcommand{\SetU}{\Set{U}}
\newcommand{\SetP}{\Set{P}}
\newcommand{\SetV}{\Set{V}}
\newcommand{\SetUi}{\Seti{U}{i}}
\newcommand{\SetT}{\Set{T}}
\newcommand{\SetTp}{\Set{T}'}
\newcommand{\SetTb}{\Set{T_\textit{b}}}
\newcommand{\SetM}{\Set{M}}

\newcommand{\SM}{\mat{S_m}}
\newcommand{\SP}{\mat{S_p}}
\newcommand{\SF}{\mat{S_f}}
\newcommand{\SW}{\mat{S_w}}
\newcommand{\IdentMat}[1]{\mat{I_{#1 \times #1}}}
\newcommand{\ZeroMat}[1]{\mat{0_{#1 \times #1}}}

% Offroad-Driving
\newcommand{\real}{\mathcal{R}}
\newcommand{\Map}{\mathcal{M}}
\newcommand{\map}{\mathbf{m}}
\newcommand{\Perception}{\mathcal{P}}